%  LaTeX support: latex@mdpi.com 
%  For support, please attach all files needed for compiling as well as the log file, and specify your operating system, LaTeX version, and LaTeX editor.

%=================================================================
\documentclass[jmse,article,submit,pdftex,moreauthors]{Definitions/mdpi} 

%=================================================================

% MDPI internal commands
\firstpage{1} 
\makeatletter 
\setcounter{page}{\@firstpage} 
\makeatother
\pubvolume{1}
\issuenum{1}
\articlenumber{0}
\pubyear{2024}
\copyrightyear{2024}
%\externaleditor{Academic Editor: Firstname Lastname}
\datereceived{30 June 2024} 
\daterevised{} 
\dateaccepted{} 
\datepublished{} 
\hreflink{https://doi.org/} % If needed use \linebreak
%\doinum{}

%=================================================================
% Add packages and commands here. The following packages are loaded in our class file: fontenc, inputenc, calc, indentfirst, fancyhdr, graphicx, epstopdf, lastpage, ifthen, lineno, float, amsmath, setspace, enumitem, mathpazo, booktabs, titlesec, etoolbox, tabto, xcolor, soul, multirow, microtype, tikz, totcount, changepage, attrib, upgreek, cleveref, amsthm, hyphenat, natbib, hyperref, footmisc, url, geometry, newfloat, caption

\graphicspath{{figures/}} % path to folder with figures
\usepackage{subcaption}
\usepackage{listings}
\usepackage{xcolor}
\usepackage{gensymb} % for \textdegree
\usepackage{tabularx}
\newcolumntype{Y}{>{\centering\arraybackslash}X}
\usepackage{amsmath,mathtools,amsthm}
	\setcounter{MaxMatrixCols}{15}
\usepackage{array}
\usepackage{amssymb}
\usepackage{amsfonts}
\usepackage{float}
\usepackage{chngcntr}
\counterwithin*{equation}{section}
\usepackage[super]{nth}
\usepackage{longtable}
\usepackage{tabularx}
\usepackage{multirow}
\usepackage{adjustbox}

\definecolor{deepblue}{rgb}{0,0,0.5}
\definecolor{deepred}{rgb}{0.6,0,0}
\definecolor{deepmagenta}{RGB}{195,12,214}
\definecolor{deepgreen}{RGB}{29,122,30}
\definecolor{mintgreen}{RGB}{192,249,192}
\definecolor{codepurple}{RGB}{204, 0, 153}
\definecolor{LightCyan}{rgb}{0.88,1,1}
\definecolor{GainsboroPurple}{RGB}{247,176,247}
\definecolor{LightSlateGrey}{RGB}{124,118,118}
%    
\lstdefinestyle{mystyle}{
    commentstyle=\color{LightSlateGrey},
    keywordstyle=\color{deepgreen}\bfseries,
    emphstyle=\color{deepred},
    numberstyle=\tiny\color{deepblue},
    stringstyle=\color{codepurple},
    basicstyle=\ttfamily\scriptsize,
    breakatwhitespace=false,         
    breaklines=true,                 
    captionpos=b,                    
    keepspaces=true,                 
    numbers=left,                    
    numbersep=5pt,                  
    showspaces=false,                
    showstringspaces=false,
    showtabs=false,                  
    tabsize=2
}
\lstset{style=mystyle}

%=================================================================
% Full title of the paper (Capitalized)
\Title{Artificial Intelligence for Computational Remote Sensing: Quantifying Patterns of Land Cover Types in Port Phillip Bay, Australia}

% MDPI internal command: Title for citation in the left column
\TitleCitation{Artificial Intelligence for Computational Remote Sensing: Quantifying Patterns of Land Cover Types in Port Phillip Bay, Australia}

% Author Orchid ID: enter ID or remove command
\newcommand{\orcidauthorA}{0000-0002-5759-1089} % Polina Add \orcidA{} behind the author's name

% Authors, for the paper (add full first names)
\Author{Polina Lemenkova $^{1}$,$^{2}$*\orcidA{}}

%\longauthorlist{yes}

% MDPI internal command: Authors, for metadata in PDF
\AuthorNames{Polina Lemenkova}

% MDPI internal command: Authors, for citation in the left column
\AuthorCitation{Lemenkova, P.}
% If this is a Chicago style journal: Lastname, Firstname, Firstname Lastname, and Firstname Lastname.

% Affiliations / Addresses (Add [1] after \address if there is only one affiliation.)
\address{%
$^{1}$ \quad Department of Geoinformatics, Faculty of Digital and Analytical Sciences, Paris Lodron Universität Salzburg, Schillerstraße 30, A-5020 Salzburg, Austria.

$^{2}$ \quad Department of Biological, Geological and Environmental Sciences, Alma Mater Studiorum -- Università di Bologna, Via Irnerio 42, Bologna, IT-40126, Emilia-Romagna, Italy.
}

% Contact information of the corresponding author
\corres{Correspondence: polina.lemenkova@plus.ac.at; polina.lemenkova2@unibo.it; Tel.: (+43)-677-6173-2772 or (+39)-344-69-28-732}

% Current address and/or shared authorship
%\firstnote{Current address: Affiliation 3.} 

% Abstract (Do not insert blank lines, i.e. \\) 
\abstract{A single paragraph of about 200 words maximum. }

% Keywords
\keyword{image processing; Earth observation; coastal management; GIS; cartography; geoinformatics; inland water quality; landscape analysis; environmental monitoring; sensors; machine learning} 

% The fields PACS, MSC, and JEL may be left empty or commented out if not applicable
%\PACS{J0101}
%\MSC{}
%\JEL{}

\begin{document}

%----------- section ----------->
\section{Introduction}

%----------- section ----------->
\section{Study Area}

The general topographic map of Australia with outlined location of the study area was plotted based on the General Bathymetric Chart of the Oceans (GEBCO) dataset with 15 arc minute resolution. The names of the major features, cities and geographic objects were obtained from the Gazetteer of Australia (Aggregator: Geoscience Australia, distributed under the Creative Commons Attribution 3.0 license). 

\begin{figure}[H]
    \begin{center}
	\hspace{-35pt}\includegraphics[width=14.0cm]{Fig_01.jpg}
    \end{center}
    \caption{General topographic map of Australia with outlined location of the study area. Mapping software: Generic Mapping Tools (GMT) version 6.4.0. Data source: GEBCO. Map source: author.}
    \label{fig01}
\end{figure}

\begin{figure}[H]
    \begin{center}
	\hspace{-35pt}\includegraphics[width=14.0cm]{Fig_01.jpg}
    \end{center}
    \caption{Enlarged fragment of the topographic map of southern Australia showing the location of the study area. Mapping software: Generic Mapping Tools (GMT) version 6.4.0. Hydrographic and topographic data source: GEBCO. Map source: author.}
    \label{fig02}
\end{figure}

The data source for the land cover types reference were derived from the nationally consistent and thematically comprehensive land cover reference for Australia -- the Dynamic Land Cover Dataset (DLCD), Figure \ref{fig03}.

\begin{figure}[H]
    \begin{center}
	\hspace{-35pt}\includegraphics[width=14.0cm]{Fig_03.jpg}
    \end{center}
    \caption{General topographic map of Australia with outlined location of the study area. Mapping software: QGIS version 3.34.7 'Prizren'. Data source: Dynamic Land Cover Dataset (DLCD). Map source: author.}
    \label{fig03}
\end{figure}

%----------- section ----------->
\section{Materials and Methods}

%----------- subsection ----------->
\subsection{Data}

\begin{figure}[H]
  \centering
  \begin{tabular}[b]{c}
    \includegraphics[width=2.9cm]{Fig_04a.jpg} \\
    \small (a): March 2013
  \end{tabular}
  \begin{tabular}[b]{c}
    \includegraphics[width=2.9cm]{Fig_04b.jpg} \\
    \small (b): March 2015
  \end{tabular}
    \begin{tabular}[b]{c}
    \includegraphics[width=2.9cm]{Fig_04c.jpg} \\
    \small (c): March 2017
  \end{tabular}
  \begin{tabular}[b]{c}
    \includegraphics[width=2.9cm]{Fig_04d.jpg} \\
    \small (d): March 2024
  \end{tabular}
  \caption{Original data: Landsat 8-9 OLI/TIRS images on southern Australia, Philippe Bay collected during the autumn period (March) on years 2013, 2015, 2017 and 2024. Source: USGS \cite{Landsat}. Compilation source: author.}
  \label{fig04}
\end{figure}

%----------- subsection ----------->
\subsection{Subsection}

%----------- subsection ----------->
\subsection{Subsection}

%----------- section ----------->
\section{Results}

\begin{figure}[H]
  \centering
  \begin{tabular}[b]{c}
    \includegraphics[width=2.9cm]{Fig_04a.jpg} \\
    \small (a): Random Forest
  \end{tabular}
  \begin{tabular}[b]{c}
    \includegraphics[width=2.9cm]{Fig_04b.jpg} \\
    \small (b): SVM
  \end{tabular}
    \begin{tabular}[b]{c}
    \includegraphics[width=2.9cm]{Fig_04c.jpg} \\
    \small (c): ANN;
  \end{tabular}
  \begin{tabular}[b]{c}
    \includegraphics[width=2.9cm]{Fig_04d.jpg} \\
    \small (d): DTC
  \end{tabular}
  \caption{Results of image processing on 2013 using four different methods: 1) Random Forest; 2) Support Vector Machine (SVM); 3) ANN-based approach using MLPClassifier; 4) DecisionTreeClassifier. Image processing source: author.}
  \label{fig05}
\end{figure}

\begin{figure}[H]
  \centering
  \begin{tabular}[b]{c}
    \includegraphics[width=2.9cm]{Fig_04a.jpg} \\
    \small (a): Random Forest
  \end{tabular}
  \begin{tabular}[b]{c}
    \includegraphics[width=2.9cm]{Fig_04b.jpg} \\
    \small (b): SVM
  \end{tabular}
    \begin{tabular}[b]{c}
    \includegraphics[width=2.9cm]{Fig_04c.jpg} \\
    \small (c): ANN
  \end{tabular}
  \begin{tabular}[b]{c}
    \includegraphics[width=2.9cm]{Fig_04d.jpg} \\
    \small (d): DTC
  \end{tabular}
  \caption{Results of image processing on 2015 using four different methods: 1) Random Forest; 2) Support Vector Machine (SVM); 3) ANN-based approach using MLPClassifier; 4) DecisionTreeClassifier. Image processing source: author.}
  \label{fig06}
\end{figure}

\begin{figure}[H]
  \centering
  \begin{tabular}[b]{c}
    \includegraphics[width=2.9cm]{Fig_04a.jpg} \\
    \small (a): Random Forest
  \end{tabular}
  \begin{tabular}[b]{c}
    \includegraphics[width=2.9cm]{Fig_04b.jpg} \\
    \small (b): SVM
  \end{tabular}
    \begin{tabular}[b]{c}
    \includegraphics[width=2.9cm]{Fig_04c.jpg} \\
    \small (c): ANN
  \end{tabular}
  \begin{tabular}[b]{c}
    \includegraphics[width=2.9cm]{Fig_04d.jpg} \\
    \small (d): DTC
  \end{tabular}
  \caption{Results of image processing on 2017 using four different methods: 1) Random Forest; 2) Support Vector Machine (SVM); 3) ANN-based approach using MLPClassifier; 4) DecisionTreeClassifier. Image processing source: author.}
  \label{fig07}
\end{figure}

\begin{figure}[H]
  \centering
  \begin{tabular}[b]{c}
    \includegraphics[width=2.9cm]{Fig_04a.jpg} \\
    \small (a): Random Forest
  \end{tabular}
  \begin{tabular}[b]{c}
    \includegraphics[width=2.9cm]{Fig_04b.jpg} \\
    \small (b): SVM
  \end{tabular}
    \begin{tabular}[b]{c}
    \includegraphics[width=2.9cm]{Fig_04c.jpg} \\
    \small (c): ANN
  \end{tabular}
  \begin{tabular}[b]{c}
    \includegraphics[width=2.9cm]{Fig_04d.jpg} \\
    \small (d): DTC
  \end{tabular}
  \caption{Results of image processing on 2024 using four different methods: 1) Random Forest; 2) Support Vector Machine (SVM); 3) ANN-based approach using MLPClassifier; 4) DecisionTreeClassifier. Image processing source: author.}
  \label{fig08}
\end{figure}

\begin{figure}[H]
  \centering
  \begin{tabular}[b]{c}
    \includegraphics[width=2.9cm]{Fig_04a.jpg} \\
    \small (a): 2013
  \end{tabular}
  \begin{tabular}[b]{c}
    \includegraphics[width=2.9cm]{Fig_04b.jpg} \\
    \small (b): 2015
  \end{tabular}
    \begin{tabular}[b]{c}
    \includegraphics[width=2.9cm]{Fig_04c.jpg} \\
    \small (c): 2017
  \end{tabular}
  \begin{tabular}[b]{c}
    \includegraphics[width=2.9cm]{Fig_04d.jpg} \\
    \small (d): 2024
  \end{tabular}
  \caption{Time series of the classified images (2013-2024) processed using k-means clustering method. Image processing source: author.}
  \label{fig09}
\end{figure}

\begin{figure}[H]
  \centering
  \begin{tabular}[b]{c}
    \includegraphics[width=2.9cm]{Fig_04a.jpg} \\
    \small (a): 2013
  \end{tabular}
  \begin{tabular}[b]{c}
    \includegraphics[width=2.9cm]{Fig_04b.jpg} \\
    \small (b): 2015
  \end{tabular}
    \begin{tabular}[b]{c}
    \includegraphics[width=2.9cm]{Fig_04c.jpg} \\
    \small (c): 2017
  \end{tabular}
  \begin{tabular}[b]{c}
    \includegraphics[width=2.9cm]{Fig_04d.jpg} \\
    \small (d): 2024
  \end{tabular}
  \caption{Accuracy analysis of image classification (2013-2024) using chi square algorithm method. Image processing source: author.}
  \label{fig10}
\end{figure}

%----------- section ----------->
\section{Discussion}

%----------- section ----------->
\section{Conclusions}

%----------- section ----------->
%\authorcontributions{Supervision, conceptualization, methodology, software, resources, funding acquisition, and project administration, writing---original draft preparation, methodology, software, data curation, visualization, formal analysis, validation, writing---review and editing, and investigation, P.L.}

\funding{The publication was funded by the Editorial Office of JMSE, Multidisciplinary Digital Publishing Institute (MDPI), by providing 100\% discount for the APC of this manuscript.}

\institutionalreview{Not applicable.}

\informedconsent{Not applicable.}

\dataavailability{Not applicable} 

\acknowledgments{The authors thank the reviewers for reading and review of this manuscript.}

\conflictsofinterest{The authors declare no conflict of interest.} 

\abbreviations{Abbreviations}{
The following abbreviations are used in this manuscript:\\

\noindent 
\begin{tabular}{@{}ll}
AI & Artificial Intelligence \\
ANN & Artificial Neural Network \\
DN & Digital Number \\
EO & Earth Observation \\
GEBCO & General Bathymetric Chart of the Oceans \\
GMT & Generic Mapping Tools \\
GRASS & Geographic Resources Analysis Support System \\
GIS & Geographic Information System \\
Landsat OLI/TIRS & Landsat Operational Land Imager and Thermal Infrared Sensor \\
ML & Machine Learning \\
MLP & Multilayer Perceptron \\
RF & Random Forest \\
RS & Remote Sensing \\
SVM & Support Vector Machine \\
USGS & United States Geological Survey \\
UTM & Universal Transverse Mercator \\
WGS84 & World Geodetic System 84 \\
WRS & Worldwide Reference System
\end{tabular}
}

%%%%%%%%%%%%%%%%%%%%%%%%%%%%%%%%%%%%%%%%%%
%% Optional
\appendixtitles{no} % Leave argument "no" if all appendix headings stay EMPTY (then no dot is printed after "Appendix A"). If the appendix sections contain a heading then change the argument to "yes".
\appendixstart
\appendix
\section[\appendixname~\thesection]{}
\subsection[\appendixname~\thesubsection]{}
The appendix 

\section[\appendixname~\thesection]{}
All appendix sections must be cited in the main text. In the appendices, Figures, Tables, etc. should be labeled, starting with ``A''---e.g., Figure A1, Figure A2, etc.

%%%%%%%%%%%%%%%%%%%%%%%%%%%%%%%%%%%%%%%%%%
\begin{adjustwidth}{-\extralength}{0cm}
%\printendnotes[custom] % Un-comment to print a list of endnotes

\reftitle{References}
\nocite{*}

%=====================================
% References, variant A: external bibliography
%=====================================
\bibliography{Australia.bib}

%%%%%%%%%%%%%%%%%%%%%%%%%%%%%%%%%%%%%%%%%%
\end{adjustwidth}
\end{document}
